\section*{Chapter \ref{ch:b2:sound-waves} --- Sound Waves in Fluids}

\begin{solution}{exo:b2:sound-waves:1}
$c = \sqrt{\gamma RT/M}$: \qty{331}{m/s}, \qty{343}{m/s}, \qty{299}{m/s}; helium
\qty{1010}{m/s}; CO$_2$ \qty{268}{m/s}. The vocal cords vibrate at the same
frequency, but the resonances of the vocal tract (which shape the
timbre) scale with $c$ and move up by a factor three: the voice sounds
high-pitched.
\end{solution}

\begin{solution}{exo:b2:sound-waves:2}
\qty{60}{dB}: $I = \qty{1e-6}{W/m^2}$, $p_m = \sqrt{2ZI} = \qty{29}{mPa}$, $v_m = p_m/Z =
\qty{7e-5}{m/s}$, $\xi_m = v_m/\omega = \qty{22}{nm}$. \qty{120}{dB}: \qty{1}{W/m^2},
\qty{29}{Pa}, \qty{7}{cm/s}, \qty{22}{\micro m}. Two: \qty{63}{dB}; a hundred: \qty{80}{dB}.
\end{solution}

\begin{solution}{exo:b2:sound-waves:3}
$I = 1/4\pi r^2$: \qty{0.080}{W/m^2} (\qty{109}{dB}), \qty{89}{dB}, \qty{69}{dB}; $I = 10^{-6}$
at $r = \sqrt{1/4\pi10^{-6}} = \qty{280}{m}$; \qty{100}{dB} at \qty{30}{m} needs $4\pi \times
900 \times 10^{-2} = \qty{110}{W}$ of sound.
\end{solution}

\begin{solution}{exo:b2:sound-waves:4}
$4.5 \times 340 = \qty{1.5}{km}$; $340 \times 1.5 = \qty{510}{m}$; $1500 \times 1.2 = \qty{1.8}{km}$;
$1540 \times 65 \times 10^{-6} = \qty{10}{cm}$.
\end{solution}

\begin{solution}{exo:b2:sound-waves:5}
(a) $500 \times 340/300 = \qty{567}{Hz}$; $500 \times 340/380 = \qty{447}{Hz}$. (b) $500 \times
380/340 = \qty{559}{Hz}$; both: $500 \times 380/300 = \qty{633}{Hz}$, against \qty{654}{Hz}
(source alone at \qty{80}{m/s}) or \qty{618}{Hz} (observer alone): the air, not
the relative speed, sets the result. (c) The wall receives \qty{567}{Hz}
and re-emits it at rest; the driver approaches at \qty{40}{m/s}: $567 \times
380/340 = \qty{633}{Hz}$.
\end{solution}

\begin{solution}{exo:b2:sound-waves:6}
$v_m = p_m/Z$: \qty{2.4}{mm/s}, \qty{0.67}{\micro m/s}, \qty{25}{nm/s}; $\xi_m = v_m/\omega$:
\qty{390}{nm}, \qty{0.11}{nm}, \qty{4}{pm}; $I = p_m^2/2Z$: \qty{1.2e-3}{W/m^2} (\qty{91}{dB}),
\qty{3.3e-7}{W/m^2} (\qty{55}{dB}), \qty{1.3e-8}{W/m^2} (\qty{41}{dB}). A stiff medium
barely moves under a given pressure and carries little power.
\qty{1}{Pa} is \qty{120}{dB} re \qty{1}{\micro Pa}.
\end{solution}

\begin{solution}{exo:b2:sound-waves:7}
(a) $\sqrt{1.013 \times 10^5/1.2} = \qty{290}{m/s}$; $\times\sqrt{1.4}$: \qty{343}{m/s}. (b)
$\gamma = c^2M/RT$: $1.40$ for air, $1.71$ for argon ($5/3$ within the rounding):
three translational degrees of freedom ($\gamma = 1 + 2/3$) against five for
a diatomic gas ($1 + 2/5$). (c) $\sqrt{D/f} = \sqrt{2 \times 10^{-8}} = \qty{0.14}{mm}
\ll \lambda = \qty{34}{cm}$: no heat flows between crests and troughs.
\end{solution}

\begin{solution}{exo:b2:sound-waves:8}
(a) Open--open: $p = p_m\sin kx\cos\omega t$ with $kL = n\pi$; closed--open:
$p = p_m\cos kx\cos\omega t$ (velocity node, pressure antinode at the closed
end) with $kL = (2n - 1)\pi/2$. (b) $L = c/2f = \qty{39}{cm}$; $c/4f = \qty{19.5}{cm}$.
(c) All harmonics; odd harmonics only (the clarinet's hollow sound).
(d) $0.6 \times 2 = \qty{1.2}{cm}$ per open end: $L_{\text{eff}} = \qty{41.4}{cm}$, $f =
\qty{414}{Hz}$, $6\%$ flat --- the pipe must be cut to \qty{36.6}{cm}.
\end{solution}

\begin{solution}{exo:b2:sound-waves:9}
$\mathcal P = \qty{400}{W}$, $I = \mathcal P/2\pi r^2$. (a) \qty{6.4e-3}{W/m^2}, \qty{98}{dB},
$p_m = \qty{2.3}{Pa}$. (b) $r = \sqrt{400/2\pi} = \qty{8}{m}$; \qty{70}{dB} at \qty{2.5}{km}.
(c) $98 - 20\log(r/100) - (r - 100)/100 = 70$: $r \approx \qty{950}{m}$. (d)
$\pi(10)^2/\pi(0.95)^2 \approx 110$ sirens.
\end{solution}

\begin{solution}{exo:b2:sound-waves:10}
(a) $\rho_0\partial_tv = p_mk\sin kx\cos\omega t$: $v = (p_m/Z)\sin kx\sin\omega t$. (b)
$e_k = \tfrac12\chi_Sp_m^2\sin^2kx\sin^2\omega t$, $e_p = \tfrac12\chi_Sp_m^2\cos^2kx\cos^2\omega t$:
one is maximal when the other vanishes, in time and in space. (c) $E =
S\int_0^L(e_k + e_p)\dd x = \tfrac14\chi_Sp_m^2SL$, constant. (d) $I = pv =
(p_m^2/4Z)\sin2kx\sin2\omega t$: zero mean, but twice per period energy
flows from the pressure antinodes to the velocity antinodes and back.
\end{solution}

\begin{solution}{exo:b2:sound-waves:11}
(a) The front emitted at $t = 0$ has radius $ct$ when the source is $vt$
away: the tangent cone has $\sin\theta = c/v$. (b) $\theta = \qty{39}{\degree}$; the cone
reaches the ground under the aircraft when it is $h/\tan\theta = \qty{15}{km}$
further: $15000/480 = \qty{31}{s}$ after the overhead passage. (c) A
\qty{60}{km}-wide carpet along the whole route hears the boom. (d) $\sin\theta
= 340/800$, $\theta = \qty{25}{\degree}$; the crack of the cone arrives first, the
muzzle blast, travelling at $c$ from the gun, later.
\end{solution}

\begin{solution}{exo:b2:sound-waves:12}
(a) $f_1 = f(c + v\cos\theta)/c$. (b) $f_2 = f_1c/(c - v\cos\theta) = f(c + v\cos\theta)
/(c - v\cos\theta) \approx f(1 + 2v\cos\theta/c)$. (c) $2 \times 5 \times 10^6 \times 0.5 \times 0.5/1540
= \qty{1.6}{kHz}$; zero at \qty{90}{\degree}: tilt the probe. (d) $\Delta v = 50 \times 1540/
(2 \times 5 \times 10^6 \times 0.5) = \qty{1.5}{cm/s}$. The ear follows pitch and timbre in
real time: a smooth whistle means \omterm{def:b2:viscous-flows:reynolds}{laminar flow}, a hiss means turbulence
behind a narrowing.
\end{solution}

\begin{solution}{pb:b2:sound-waves:1}
\textbf{1.} $c = \sqrt{1.4 \times 8.31 \times 288/0.029} = \qty{340}{m/s}$; \qty{349}{m/s}
at \qty{30}{\degreeCelsius}. \qty{1}{km} in \qty{2.9}{s}.

\textbf{2.} \qty{1.4}{km}. The channel is kilometres long; the sound of
its different parts arrives over seconds, with echoes from clouds and
ground.

\textbf{3.} $I = \qty{1}{W/m^2}$: $p_m = \sqrt{2ZI} = \qty{29}{Pa}$, $v_m = \qty{7}{cm/s}$.

\textbf{4.} $It = \qty{2}{J/m^2}$.

\textbf{5.} $c$ grows upward; $\sin i/c$ constant makes $i$ grow, so the
ray bends toward the horizontal and back down: sound is channelled
along the ground and carries far. Over water the lowest air is cooled
by the water: the same.

\textbf{6.} The rays bend upward, leaving a shadow zone. Radius of
curvature $R = c/(\dd c/\dd z) = \qty{50}{km}$: after \qty{3}{km} the ray has
turned by $3/50 = \qty{0.06}{rad} = \qty{3.4}{\degree}$ and risen \qty{90}{m}: the
listener is below it.

\textbf{7.} Absorption grows as $f^2$: the crack's high frequencies are
gone after a few kilometres, the low rumble remains.

\textbf{8.} $\mathcal P = \qty{2}{W}$; $I = 2/2\pi r^2 = \qty{3.2e-3}{W/m^2}$ at \qty{10}{m}:
\qty{95}{dB}.

\textbf{9.} $p_m = \sqrt{2 \times 410 \times 3.2 \times 10^{-3}} = \qty{1.6}{Pa}$; $v_m = \qty{3.9}{mm/s}$;
$\xi_m = v_m/2\pi f = \qty{6}{\micro m}$.

\textbf{10.} $+\qty{20}{dB}$: \qty{115}{dB} at \qty{1}{m} --- damaging within
minutes. \qty{60}{dB}: $r = \sqrt{2/2\pi10^{-6}} = \qty{560}{m}$.

\textbf{11.} Ten incoherent: $+\qty{10}{dB}$, \qty{105}{dB}. Two in phase: the
pressure doubles, $+\qty{6}{dB}$, \qty{101}{dB}.

\textbf{12.} Balance $\mathcal P = \alpha A\,ce/4$: $e = 8/(0.3 \times 2000 \times 340) =
\qty{3.9e-5}{J/m^3}$, $E = eV = \qty{0.4}{J}$. After the stop, $\dd E/\dd t = -(\alpha Ac
/4V)E$: $\tau = 4V/\alpha Ac = \qty{0.20}{s}$, and \qty{60}{dB} takes $\tau\ln10^6 =
\qty{2.7}{s}$ (Sabine's reverberation time).

\textbf{13.} $\lambda = c/f$: \qty{0.51}{mm}, \qty{0.31}{mm}, \qty{0.15}{mm}.

\textbf{14.} $2 \times 0.12/1540 = \qty{156}{\micro s}$; \qty{6.4}{kHz}; $200$ lines in
\qty{31}{ms} (thirty images a second).

\textbf{15.} Tissue--air: $r = -0.9995$, $99.9\%$ of the energy reflected
--- hence the gel, which removes the air film. Tissue--bone: $r = 0.63$,
$40\%$ of the energy; the rest is absorbed by the bone, which casts a
shadow. Fat--muscle: $r = 0.10$, $1\%$ of the energy: the faint echoes the
image is made of.

\textbf{16.} Reflected fraction $1.1\%$ ($-\qty{20}{dB}$); round trip
\qty{4}{cm} at \qty{5}{MHz}: $-\qty{10}{dB}$; in all $-\qty{30}{dB}$, one thousandth.

\textbf{17.} Round trip \qty{24}{cm}: \qty{36}{dB}, \qty{60}{dB}, \qty{120}{dB}: \qty{3}{MHz}
and \qty{5}{MHz} work, \qty{10}{MHz} does not --- deep organs are imaged at
low frequency with coarse resolution, superficial ones at high
frequency finely.

\textbf{18.} Five wavelengths, \qty{1.5}{mm} long: two interfaces closer
than about \qty{0.8}{mm} (half the pulse) give overlapping echoes.

\textbf{19.} $2 \times 5 \times 10^6 \times 0.4 \times 0.5/1540 = \qty{1.3}{kHz}$: a whistle whose
pitch rises and falls with each heartbeat.

\textbf{20.} $v = \qty{25}{m/s}$: $700 \times 340/315 = \qty{756}{Hz}$, $700 \times 340/365 =
\qty{652}{Hz}$; ratio $1.16$, $12\log_21.16 = 2.6$ semitones.

\textbf{21.} $700 \times 355/315 = \qty{789}{Hz}$.

\textbf{22.} At closest approach, when the ambulance's velocity is
perpendicular to the line of sight (strictly, when the sound now
arriving was emitted at that point).

\textbf{23.} $\lambda' = (340 - 25)/700 = \qty{0.45}{m}$ (against \qty{0.49}{m}); the
sound recedes from the ambulance at \qty{315}{m/s} ahead and \qty{365}{m/s}
behind.

\textbf{24.} $v = c\,\Delta f/2f = 3 \times 10^8 \times 3200/4.8 \times 10^{10} = \qty{20}{m/s} =
\qty{72}{km/h}$.

\textbf{25.} \qty{340}{m/s}, \qty{1540}{m/s}, \qty{25}{m/s}, \qty{20}{m/s}, \qty{0.4}{m/s},
\qty{3e8}{m/s}: $f' = f(c + v_o)/(c - v_s)$, applied once or twice, and for
light at low speed.
\end{solution}
